% Research Document: Scale Diffusion for Molecular Generation
% & Multifunctional Coating Generative AI Platform
% Compile with: pdflatex research_document.tex
\documentclass[11pt,a4paper]{article}

\usepackage[T1]{fontenc}
\usepackage{lmodern}
\usepackage[margin=2.5cm]{geometry}
\usepackage{hyperref}
\usepackage{amsmath,amssymb}
\usepackage{graphicx}
\usepackage{booktabs}
\usepackage{enumitem}
\usepackage{xcolor}
\usepackage[numbers,sort&compress]{natbib}

\definecolor{secblue}{HTML}{1971C2}
\definecolor{linkblue}{HTML}{4263EB}

\hypersetup{
  colorlinks=true,
  linkcolor=secblue,
  citecolor=secblue,
  urlcolor=linkblue,
}

\title{\textbf{Scale Diffusion for Molecular Generation \&\\
Multifunctional Coating Generative AI Platform}\\[6pt]
\large Comprehensive Research Survey}
\author{MolSSD Project}
\date{March 2026}

\begin{document}
\maketitle
\tableofcontents
\newpage

%% ========================================================================
\section{Executive Summary}
%% ========================================================================

This document surveys the scientific landscape underpinning the
\textbf{MolSSD} (Molecular Scale Space Diffusion) project: a generative AI
platform that applies scale-space diffusion theory to molecular generation,
targeting the inverse design of novel multifunctional coating materials.
The platform bridges computer vision (scale-space theory), statistical
mechanics (coarse-graining), equivariant deep learning (E(n)-GNNs), and
materials chemistry (UV absorbers, anti-corrosion, self-healing coatings).

The key innovation is adapting the recently published Scale Space Diffusion
(SSD) framework---originally developed for image generation---to molecular
graphs, creating a principled hierarchical generative model that operates
across multiple resolutions: from single centroids through chromophore
skeletons and molecular fragments down to full atomic detail.

%% ========================================================================
\section{Scale-Space Diffusion Theory}
%% ========================================================================

\subsection{Classical Scale-Space Theory}

Scale-space theory, originating in computer vision with Witkin (1983),
Koenderink (1984), and Lindeberg (1994), provides a mathematical framework
for representing signals at multiple levels of detail. The core idea is that
a signal can be embedded in a one-parameter family of progressively smoothed
(coarsened) versions, where fine-scale details are successively suppressed
while coarse-scale structure is preserved.

In the continuous setting, the linear scale-space of an image $f(x)$ is
generated by convolution with Gaussian kernels of increasing width:
\begin{equation}
  L(x; t) = G(\cdot; t) * f(x), \quad
  G(x; t) = \frac{1}{(2\pi t)^{d/2}} \exp\!\Bigl(-\frac{\|x\|^2}{2t}\Bigr)
\end{equation}
where $t$ is the scale parameter. This process satisfies several axiomatic
properties: causality (no new features are created at coarser scales),
semi-group structure, and translational/rotational invariance.

\subsection{Scale Space Diffusion (SSD)}

The Scale Space Diffusion framework (arXiv:2603.08709, March 2026) unifies
classical scale-space theory with denoising diffusion probabilistic models
(DDPMs). Rather than using only isotropic Gaussian noise injection (as in
standard DDPM), SSD composes \emph{resolution-reducing blur operators} with
noise:

\begin{equation}
  x_t = M_{1:t}\, x_0 + \sigma_t\, \varepsilon, \quad
  \varepsilon \sim \mathcal{N}(0, I)
\end{equation}

where $M_{1:t} = M_t \cdots M_2 M_1$ is a composed degradation operator that
progressively blurs (coarsens) the data. The key insight is that sampling
begins at low resolution and progressively adds detail, achieving a
\textbf{1.58$\times$ inference speedup} compared to standard DDPM for
$64 \times 64$ image generation while maintaining quality.

A recent analysis establishes that a wide variety of existing probabilistic
diffusion models fulfil generalised scale-space properties, creating a new
class of scale-spaces that considers the evolution of \emph{probability
distributions} rather than deterministic signals.

\subsection{Non-Isotropic Posteriors}

At resolution-changing steps (where $M_t$ reduces dimensionality), the
reverse-process posterior becomes \emph{non-isotropic}:
\begin{equation}
  \Sigma = \sigma_{t-1}^2 I - \frac{\sigma_{t-1}^4}{\sigma_t^2}\, M_t^\top M_t
\end{equation}

Computing this covariance requires eigendecomposition of $M_t^\top M_t$.
The \textbf{Lanczos algorithm} provides an efficient solution: it computes
the required Ritz pairs via matrix-vector products (never forming the full
matrix), with $\mathcal{O}(Nk^2)$ complexity where $k$ is the number of
super-nodes---far cheaper than the $\mathcal{O}(N^3)$ dense
eigendecomposition.

\subsection{Rate-Distortion Interpretation}

SSD admits an information-theoretic interpretation through rate-distortion
theory. Each scale level corresponds to a particular rate-distortion
trade-off: coarser representations carry fewer bits but introduce more
distortion. This provides a principled basis for choosing the number of
scale levels and the allocation of denoising steps per resolution, filling
a gap where previous multi-scale molecular methods relied on ad-hoc
choices.

%% ========================================================================
\section{Molecular Scale Space}
%% ========================================================================

\subsection{Spectral Graph Coarsening}

The central technical challenge in adapting SSD to molecules is replacing
Gaussian blurring (which acts on regular pixel grids) with a mathematically
principled coarsening operator for molecular graphs. MolSSD uses
\textbf{spectral graph coarsening}:

\begin{enumerate}[leftmargin=2em]
  \item \textbf{Graph Laplacian}: Compute the normalised Laplacian
    $L = I - D^{-1/2}AD^{-1/2}$ from the molecular adjacency matrix $A$.
  \item \textbf{Eigendecomposition}: Obtain eigenvalues $\lambda_k$ and
    eigenvectors $\phi_k$ via \texttt{torch.linalg.eigh}.
  \item \textbf{Spectral clustering}: Apply $k$-means to the lowest
    eigenvectors to partition atoms into super-node groups.
  \item \textbf{Coarsening matrices}: Construct $C^{(k)}$ mapping fine-level
    signals to super-node groups via mass-weighted aggregation.
\end{enumerate}

The low-frequency Laplacian eigenvectors naturally capture chemically
meaningful structure: the lowest frequencies give connected components,
intermediate frequencies resolve functional groups and rings, and high
frequencies encode individual atomic detail.

\subsection{Scale Hierarchy}

The molecular scale space defines a hierarchy of representations:
\begin{center}
\begin{tabular}{lll}
\toprule
\textbf{Level} & \textbf{Representation} & \textbf{Chemical Meaning} \\
\midrule
0 & Full atoms & All atomic coordinates \\
1 & Functional groups & --OH, --COOH, --NH$_2$, aromatics \\
2 & Molecular fragments & BRICS-like substructures \\
3 & Chromophore skeleton & Core conjugated system \\
4 & Single centroid & Centre of mass \\
\bottomrule
\end{tabular}
\end{center}

Unlike heuristic fragmentation methods (BRICS, Murcko scaffolds), spectral
coarsening provides a \emph{data-driven}, continuous hierarchy that adapts to
molecular topology. Recent work on \textbf{CoarsenConf} (JCIM, 2024) uses
SE(3)-equivariant hierarchical variational autoencoders with coarse-graining
based on torsional angles, while \textbf{Equivariant Blurring Diffusion}
(NeurIPS 2024) defines a forward process that moves towards fragment-level
coarse-grained structure.

\subsection{SE(3) Equivariance Preservation}

A critical requirement is that the coarsening operators preserve SE(3)
equivariance---rotations and translations of the input molecule must
produce correspondingly transformed coarsened representations. This is
achieved through \textbf{centre-of-mass (COM) aggregation}: super-node
positions are computed as mass-weighted averages of constituent atom
positions, which commutes with rigid-body transformations.

%% ========================================================================
\section{Neural Architecture}
%% ========================================================================

\subsection{Molecular Flexi-Net}

MolSSD introduces the \textbf{Molecular Flexi-Net}, a hierarchical
E(n)-equivariant graph neural network designed to operate across multiple
resolutions simultaneously:

\begin{itemize}[leftmargin=2em]
  \item \textbf{EGNN message-passing blocks}: Each layer updates both node
    features $h_i$ and coordinates $x_i$ via radial basis functions on
    interatomic distances, maintaining E(n) equivariance throughout.
  \item \textbf{Cross-scale skip connections}: Feature projections between
    coarse and fine levels enable information flow across resolutions.
  \item \textbf{Adaptive Layer Normalisation (adaLN)}: Inspired by the
    Diffusion Transformer (DiT; Peebles \& Xie, 2023), the network
    conditions on both timestep and resolution level through learnable
    scale/shift parameters.
  \item \textbf{Dynamic activation functions}: Learnable per-layer
    activations adapt to the characteristics of each resolution.
\end{itemize}

\subsection{E(n)-Equivariant Graph Neural Networks}

The EGNN architecture (Satorras et al., 2021) forms the backbone of
Flexi-Net. Recent extensions include:

\begin{itemize}[leftmargin=2em]
  \item \textbf{ViSNet} (2024): Extracts geometric features via
    vector-scalar interactive message passing, outperforming prior methods
    on MD17/MD22 benchmarks.
  \item \textbf{VN-EGNN} (2025): Integrates virtual nodes into E(n)/SE(n)-
    equivariant GNNs, setting state-of-the-art for protein binding site
    identification.
  \item \textbf{EnviroDetaNet} (2025): E(3)-equivariant message-passing
    network with multi-level representations for molecular spectra
    prediction.
  \item \textbf{ENINet} (2025): Explicitly integrates many-body equivariant
    interactions for improved quantum chemical property prediction.
\end{itemize}

\subsection{Conditioning Mechanisms}

Two conditioning strategies are employed:
\begin{enumerate}[leftmargin=2em]
  \item \textbf{Adaptive LayerNorm}: Time and resolution embeddings
    modulate hidden features via learned affine transformations at each
    layer, following the DiT paradigm.
  \item \textbf{Classifier-free guidance}: During training, conditioning
    signals (target UV spectrum, photostability constraints) are randomly
    dropped with probability $p_\text{uncond}$. At inference, guided and
    unguided predictions are interpolated:
    \begin{equation}
      \hat{\varepsilon}_\theta = (1 + w)\,\varepsilon_\theta(x_t, c) -
        w\,\varepsilon_\theta(x_t, \varnothing)
    \end{equation}
    where $w$ is the guidance weight and $c$ the conditioning signal.
\end{enumerate}

\subsection{Pre-trained Foundation Models: eSEN and UMA}

The project leverages Meta's molecular foundation models for pre-training
and transfer learning:

\textbf{eSEN} (Equivariant Scalable Energy Network) adopts a
transformer-style architecture with equivariant spherical-harmonic
representations. It improves the smoothness of the potential energy surface,
making molecular dynamics and geometry optimisations better-behaved.

\textbf{UMA} (Universal Models for Atoms) represents the latest evolution,
trained on nearly 500M unique 3D atomic structures across the OMol25, OC20,
ODAC23, OMat24, and OMC25 datasets. UMA introduces a novel
\textbf{Mixture of Linear Experts (MoLE)} architecture that enables one
model to learn from dissimilar datasets without increasing inference cost.
Both models exceed prior state-of-the-art neural network potential
performance and match high-accuracy DFT results.

%% ========================================================================
\section{Multifunctional Coating Design}
%% ========================================================================

\subsection{UV Absorbers: The Primary Target}

UV radiation causes photodegradation of polymers and coatings, threatening
the longevity of autonomous systems operating outdoors. The UV absorber
market for coatings is projected at \$1.41 billion in 2025. Current
commercial UV absorbers fall into two dominant classes:

\begin{itemize}[leftmargin=2em]
  \item \textbf{Benzotriazoles (BZTs)}: Strong absorption in the UVA range
    (300--385\,nm). Protection relies on intramolecular hydrogen bonding for
    excited-state deactivation.
  \item \textbf{Hydroxybenzophenones (HBPs)}: Moderate absorption in the
    UVB range (280--340\,nm). Function via excited-state intramolecular
    proton transfer (ESIPT).
\end{itemize}

\textbf{Limitations of existing absorbers}:
\begin{enumerate}[leftmargin=2em]
  \item Narrow absorption windows (cover only portions of UVA or UVB)
  \item Susceptibility to photodegradation under prolonged exposure
  \item Migration and leaching from polymer matrices
  \item Limited tunability outside established scaffolds
\end{enumerate}

MolSSD targets the generation of novel UV absorbers with \textbf{broad-spectrum coverage} (280--400\,nm), \textbf{exceptional photostability}
($>$1000\,h accelerated ageing), and \textbf{coating compatibility}
(solubility in polymer matrices). Theory-guided hybrid designs combining
BZT and HBP functional groups in single molecules have shown promise for
dual UVA/UVB absorption.

\subsection{Anti-Corrosion Layers}

Beyond UV protection, multifunctional coatings incorporate corrosion
inhibition:

\begin{itemize}[leftmargin=2em]
  \item \textbf{Barrier films}: Dense crosslinked polymer matrices that
    physically block moisture and ion ingress.
  \item \textbf{Inhibitor release}: Smart coatings that release corrosion
    inhibitors (e.g., cerium salts, benzotriazole) in response to pH
    changes at corrosion sites. Recent 2025 work on \textbf{pH-responsive
    self-healing coatings} uses hydrogel microdomains that undergo
    solid--liquid transitions to autonomously seal defects.
\end{itemize}

\subsection{Self-Healing Coatings}

Autonomous self-healing mechanisms represent a frontier in coating
technology:

\begin{itemize}[leftmargin=2em]
  \item \textbf{Microcapsule systems}: Healing agents (monomers, catalysts)
    encapsulated in polymer shells rupture upon mechanical damage, releasing
    contents that polymerise and seal the crack.
  \item \textbf{Intrinsic self-healing}: Reversible bonding (Diels--Alder,
    disulfide exchange, hydrogen bonding) enables the polymer matrix itself
    to reform broken bonds, sometimes triggered by heat or light.
  \item \textbf{Deep eutectic solvent (DES) approach}: A 2026 study
    demonstrated autonomous anti-corrosion coatings using hydrophobic DES
    combined with elastic resin, exhibiting scalability without performance
    degradation.
\end{itemize}

\subsection{Coatings for Autonomous Systems}

Autonomous vehicles, drones, and robotic platforms impose unique coating
requirements:
\begin{itemize}[leftmargin=2em]
  \item UV durability for sensor windows and LiDAR housings
  \item Self-cleaning hydrophobic surfaces (contact angle $>150°$)
  \item Transparency at sensor operational wavelengths
  \item Lightweight construction with environmental compliance (waterborne,
    low-VOC formulations)
\end{itemize}

%% ========================================================================
\section{Generative AI Platform Architecture}
%% ========================================================================

\subsection{Conditional Molecular Generation}

The MolSSD platform generates molecules conditioned on target properties:
\begin{itemize}[leftmargin=2em]
  \item \textbf{Spectra targets}: UV--Vis absorption profiles specified as
    target spectral envelopes; the model learns to generate molecules whose
    predicted absorption matches the target.
  \item \textbf{Multi-property guidance}: Simultaneous conditioning on
    photostability, synthetic accessibility, solubility, and coating
    compatibility through multi-head classifier-free guidance.
\end{itemize}

Recent advances in conditional molecular generation include
\textbf{DiffGui} (Nature Comms, 2025), a target-conditioned E(3)-equivariant
diffusion model integrating bond diffusion and property guidance, and
\textbf{DiffMC-Gen} (Adv.\ Sci., 2025), a dual denoising diffusion model
for multi-conditional generation.

\subsection{Closed-Loop Optimisation}

The platform implements a closed-loop workflow:
\begin{enumerate}[leftmargin=2em]
  \item \textbf{Generate}: MolSSD proposes candidate molecules conditioned on
    target properties.
  \item \textbf{Screen}: eSEN/UMA predict ground-state energies; ML
    surrogates predict UV absorption; TD-DFT validates top candidates.
  \item \textbf{Score}: Multi-objective scoring ranks candidates by property
    match, synthesisability, and novelty.
  \item \textbf{Update}: Bayesian optimisation and active learning update the
    conditioning targets and model parameters.
  \item \textbf{Repeat}: The cycle iterates until convergence or synthesis
    candidates are identified.
\end{enumerate}

This approach mirrors emerging AI-accelerated materials discovery paradigms.
AI-driven methods reduce the timeline from 10--20 years (traditional
trial-and-error) to 1--2 years through computational prediction, inverse
design, and automated experimentation.

\subsection{Multi-Objective Optimisation}

Coating molecules must simultaneously satisfy competing objectives:
\begin{itemize}[leftmargin=2em]
  \item Broad UV absorption $\leftrightarrow$ Optical transparency in visible range
  \item High photostability $\leftrightarrow$ Synthetic accessibility
  \item Strong polymer interaction $\leftrightarrow$ Low migration tendency
\end{itemize}

Pareto-front exploration via multi-objective Bayesian optimisation
identifies the best trade-off solutions, with active learning directing
computational resources to the most informative candidates.

\subsection{Self-Driving Laboratory Integration}

The platform is designed for eventual integration with autonomous
experimental systems:

\begin{itemize}[leftmargin=2em]
  \item \textbf{Polybot} (Argonne, 2026): Conducted $>$6000 battery chemical
    experiments in five months, simultaneously optimising conductivity and
    coating defects.
  \item \textbf{Rainbow}: Multi-robot SDL that autonomously optimises
    nanocrystal optical performance through closed-loop experimentation.
  \item \textbf{MINERVA}: SDL for synthesis, purification, and
    characterisation of advanced materials.
\end{itemize}

These platforms establish that AI-generated molecular candidates can be
synthesised, characterised, and fed back into the generative model in
fully autonomous cycles.

%% ========================================================================
\section{Data and Training Strategy}
%% ========================================================================

\subsection{OMol25 Pre-training}

The \textbf{Open Molecules 2025 (OMol25)} dataset from Meta FAIR provides
the foundation for pre-training:

\begin{itemize}[leftmargin=2em]
  \item $>$100 million DFT calculations at the $\omega$B97M-V/def2-TZVPD
    level of theory
  \item $\sim$83 million unique molecular systems covering small molecules,
    biomolecules, metal complexes, and electrolytes
  \item Systems of up to 350 atoms
  \item $>$6 billion CPU-hours of compute
  \item Available via HuggingFace and the fairchem library
\end{itemize}

Pre-training on OMol25 provides the model with a broad understanding of
molecular geometry and energetics before fine-tuning on task-specific data.

\subsection{Benchmark Datasets}

\begin{itemize}[leftmargin=2em]
  \item \textbf{QM9}: 134K small molecules ($\leq$9 heavy atoms) with
    13 quantum-chemical properties. Primary validation benchmark for
    comparing against EDM, GeoLDM, and other baselines.
  \item \textbf{GEOM-Drugs}: 304K drug-like conformations for scaling
    experiments. Tests the model's ability to handle larger, more flexible
    molecules with multiple conformers.
  \item \textbf{ChEMBL-3D}: $>$250 million molecular geometries for 1.8M
    drug-like compounds, optimised using AIMNet2 neural network potentials.
\end{itemize}

\subsection{Training Strategy}

\begin{itemize}[leftmargin=2em]
  \item \textbf{Min-SNR-5 loss weighting}: $w(t) =
    \min(\text{SNR}(t), \gamma) / \text{SNR}(t)$ with $\gamma=5$ for stable
    training across all timesteps.
  \item \textbf{ConvexDecay(0.5) resolution schedule}: Allocates more
    denoising steps to higher-resolution stages where atomic detail matters
    most.
  \item \textbf{EMA}: Exponential moving average with decay 0.9999 for stable
    generation quality.
  \item \textbf{Mixed precision}: bf16 training for computational efficiency.
  \item \textbf{Gradient clipping}: Max norm 1.0 to prevent training
    instabilities.
\end{itemize}

%% ========================================================================
\section{Validation Pipeline}
%% ========================================================================

\subsection{Computational Screening}

\subsubsection{TD-DFT for Excited States}

Time-dependent density functional theory (TD-DFT) at the
CAM-B3LYP/6-31+G(d,p) level provides the gold standard for predicting UV
absorption spectra. Recent work demonstrates that:

\begin{itemize}[leftmargin=2em]
  \item TD-DFT accurately reproduces experimental UV spectra of
    benzophenone derivatives and related UV absorbers.
  \item Hybrid ML/TD-DFT models (e.g., $\Delta$ML-ZINDO) achieve
    near-TDDFT accuracy at $\sim$2\,ms per molecule, enabling screening of
    millions of candidates.
  \item Interpretable ML integrated with TD-DFT descriptors and SHAP
    analysis can predict $\lambda_\text{max}$ of azo dyes with $R^2 = 0.83$.
\end{itemize}

\subsubsection{ML Surrogates}

Machine-learned surrogates provide rapid ($10^3$--$10^6\times$ faster)
property predictions:

\begin{itemize}[leftmargin=2em]
  \item \textbf{eSEN/UMA}: Ground-state energy, forces, and relaxed geometries.
  \item \textbf{Calibrated xTB}: Semi-empirical methods calibrated with ML
    give much higher accuracy than using ML to directly predict excited-state
    energies, enabling high-throughput screening.
  \item \textbf{Graph-based property predictors}: GCNs and MPNNs predict
    polymer mechanical properties, barrier efficiency, and chemical
    stability.
\end{itemize}

\subsection{Synthesisability Analysis}

\begin{itemize}[leftmargin=2em]
  \item \textbf{SA Score} (Ertl \& Schuffenhauer, 2009): Quantitative
    synthetic accessibility score based on fragment contributions. Candidates
    with SA $> 6$ are flagged as synthetically challenging.
  \item \textbf{Retrosynthesis planning}: Computer-assisted retrosynthesis
    (CASP) tools identify feasible synthetic routes from commercially
    available starting materials.
  \item \textbf{QED score}: Quantitative estimate of drug-likeness, adapted
    for materials applications as a general molecular quality metric.
\end{itemize}

\subsection{Experimental Verification}

The final validation stage involves wet-lab synthesis and characterisation:

\begin{itemize}[leftmargin=2em]
  \item \textbf{UV--Vis absorption spectroscopy}: Direct measurement of
    absorption spectra to validate computational predictions.
  \item \textbf{QUV accelerated weathering}: Standardised test (ASTM G154)
    exposing coated panels to cycles of UV radiation and moisture to assess
    photostability over $>$1000\,h equivalent outdoor exposure.
  \item \textbf{Coating formulation testing}: Solubility in target polymer
    matrices, migration resistance, and mechanical properties of the
    formulated coating.
\end{itemize}

%% ========================================================================
\section{Related Work and State of the Art}
%% ========================================================================

\subsection{Molecular Diffusion Models}

The landscape of diffusion-based molecular generation has expanded rapidly:

\begin{itemize}[leftmargin=2em]
  \item \textbf{EDM} (Hoogeboom et al., 2022): First E(n)-equivariant
    diffusion model for 3D molecules. Baseline for QM9.
  \item \textbf{GeoLDM} (Xu et al., 2023): Latent diffusion with geometric
    priors, improving generation quality through compressed representations.
  \item \textbf{HierDiff} (Qiang et al., 2023): Coarse-to-fine hierarchical
    diffusion, but relies on heuristic fragmentation rather than spectral
    coarsening.
  \item \textbf{MD3MD} (Science Advances, 2025): Multiscale graph equivariant
    diffusion that partitions molecular conformations into multiscale graphs.
  \item \textbf{EAD} (March 2026): Equivariant Asynchronous Diffusion with
    adaptive denoising schedules, achieving state-of-the-art 3D generation.
  \item \textbf{LGDC}: Latent Graph Diffusion via Spectrum-Preserving
    Coarsening---the closest existing work to MolSSD's spectral approach,
    but applied to general graphs rather than 3D molecules.
\end{itemize}

\subsection{Key Differentiators of MolSSD}

MolSSD addresses seven research gaps not covered by existing methods:

\begin{enumerate}[leftmargin=2em]
  \item \textbf{Information-theoretic basis}: SSD's rate-distortion
    interpretation provides principled scale-level allocation, unlike
    ad-hoc choices in HierDiff.
  \item \textbf{Spectral coarsening}: Data-driven hierarchy replaces
    heuristic BRICS/Murcko fragmentation.
  \item \textbf{Non-isotropic posteriors}: Formally derived and efficiently
    computed via Lanczos, vs.\ isotropic assumptions in all prior work.
  \item \textbf{Computational efficiency}: Multi-resolution denoising
    allocates more steps to fine-grained stages, reducing total compute.
  \item \textbf{Interpretable intermediates}: Scale-space levels correspond
    to chemically meaningful structures.
  \item \textbf{Unified theory}: Bridges physical coarse-graining
    (Noid et al., 2008) and generative modelling under one framework.
  \item \textbf{Application focus}: First generative design framework
    specifically targeting UV-absorbing materials for protective coatings.
\end{enumerate}

%% ========================================================================
\section{Platform Workflow Summary}
%% ========================================================================

The complete MolSSD platform operates as follows:

\begin{enumerate}[leftmargin=2em]
  \item \textbf{Pre-train} the Flexi-Net backbone on OMol25 (83M molecules,
    140M DFT calculations) to learn general molecular geometry and
    energetics.
  \item \textbf{Fine-tune} on UV absorber datasets with conditional
    generation objectives (target UV spectra, photostability, coating
    compatibility).
  \item \textbf{Generate} 10,000+ candidate molecules via scale-space
    reverse diffusion with classifier-free guidance.
  \item \textbf{Screen} candidates through a multi-stage funnel:
    \begin{itemize}
      \item Stage 1: Validity, uniqueness, novelty filters
      \item Stage 2: eSEN/UMA ground-state energy predictions
      \item Stage 3: ML surrogate UV absorption predictions
      \item Stage 4: TD-DFT validation of top-100 candidates
      \item Stage 5: SA score and retrosynthesis feasibility
    \end{itemize}
  \item \textbf{Synthesise} top 10--20 candidates in the wet lab.
  \item \textbf{Validate} via UV--Vis spectroscopy and QUV accelerated
    weathering tests.
  \item \textbf{Iterate}: Feed experimental results back into the model
    via active learning for the next generation cycle.
\end{enumerate}

%% ========================================================================
\section{Key References and Resources}
%% ========================================================================

\subsection{Core Papers}

\begin{itemize}[leftmargin=2em]
  \item \textbf{SSD}: Scale Space Diffusion (arXiv:2603.08709, March 2026)
    --- Core mathematical framework.
  \item \textbf{DiT}: Scalable Diffusion Models with Transformers
    (arXiv:2212.09748) --- Transformer backbone and adaLN conditioning.
  \item \textbf{OMol25}: Open Molecules 2025 (arXiv:2505.08762) --- Dataset
    and eSEN/UMA models.
  \item \textbf{EDM}: Equivariant Diffusion for Molecule Generation in 3D
    (arXiv:2203.17003) --- Baseline equivariant diffusion.
  \item \textbf{EGNN}: E(n) Equivariant Graph Neural Networks
    (arXiv:2102.09844) --- Core GNN architecture.
\end{itemize}

\subsection{Code Repositories}

\begin{itemize}[leftmargin=2em]
  \item SSD: \url{https://github.com/prateksha/ScaleSpaceDiffusion}
  \item EDM: \url{https://github.com/ehoogeboom/e3_diffusion_for_molecules}
  \item FairChem/eSEN: \url{https://github.com/facebookresearch/fairchem}
  \item OMol25 models: \url{https://huggingface.co/facebook/OMol25}
\end{itemize}

\subsection{Key Online Resources}

\begin{itemize}[leftmargin=2em]
  \item Flow Matching and Diffusion Models course (MIT, 2026):
    \url{https://diffusion.csail.mit.edu/2026/index.html}
  \item OMol25 documentation:
    \url{https://fair-chem.github.io/molecules/datasets/omol25.html}
  \item Pretrained models:
    \url{https://fair-chem.github.io/molecules/models.html}
\end{itemize}

%% ========================================================================
\section{Conclusion}
%% ========================================================================

The MolSSD platform represents a convergence of four research frontiers:
scale-space theory from computer vision, equivariant deep learning from
geometric ML, spectral graph theory from applied mathematics, and
inverse molecular design from computational chemistry. By adapting the
Scale Space Diffusion framework to molecular graphs, the platform provides
the first information-theoretically grounded hierarchical molecular
generation method, with direct application to the industrially significant
problem of designing novel multifunctional coating materials.

The generative AI platform extends beyond UV absorber design to encompass
anti-corrosion, self-healing, and autonomous-systems-compatible coatings,
with a closed-loop optimisation architecture that integrates computational
screening, ML surrogates, quantum-chemical validation, and experimental
feedback. This positions MolSSD as both a scientific contribution
(unifying coarse-graining and diffusion theory) and a practical tool for
accelerated materials discovery.

\end{document}
